\label{sec:result_placeholders}
This section defines the result slots for the frozen VSM run. The cells are intentionally left as dashes until all selected systems finish inference, transcript-level judging, scoring, and human audit. No table in this draft contains evidence from partial runs.

\paragraph{Main comparison.}
Table~\ref{tab:main_results} specifies the main comparison across \trace, ablations, and baselines. The final paper will state whether \tracefull\ improves over the strongest complete baseline only after the full score summary is generated.

\begin{table*}[t]
\centering
\small
\setlength{\tabcolsep}{4pt}
\begin{tabular}{lrrrrrrrr}
\toprule
System & Final & Safety & Tech. & Alliance & Action & Group & Fail. & Lat. \\
\midrule
\tracefull & -- & -- & -- & -- & -- & -- & -- & -- \\
\textsc{TRACE-NoPeer} & -- & -- & -- & -- & -- & -- & -- & -- \\
\textsc{TRACE-NoValidator} & -- & -- & -- & -- & -- & -- & -- & -- \\
\textsc{TRACE-NoSafetyCritic} & -- & -- & -- & -- & -- & -- & -- & -- \\
Stage prompt & -- & -- & -- & -- & -- & -- & -- & -- \\
Plain prompt & -- & -- & -- & -- & -- & -- & -- & -- \\
Base model & -- & -- & -- & -- & -- & -- & -- & -- \\
Prompt 1-1 & -- & -- & -- & -- & -- & -- & -- & -- \\
\bottomrule
\end{tabular}
\caption{Main VSM result slot for the frozen \texttt{vsm\_session\_core} run. Dashes mark values to be filled only after full inference, scoring, LLM judging, and human audit. ``Fail.'' reports reliability failure rate; ``Lat.'' reports average latency in seconds.}
\label{tab:main_results}
\end{table*}


\paragraph{Ablation deltas.}
Table~\ref{tab:ablation_results} specifies mechanism-level deltas relative to \tracefull. Removing peer agents should primarily affect group dynamics, removing the validator should affect stage or technique fidelity, and removing the safety critic should affect safety or boundary robustness. If an expected delta is absent, the corresponding mechanism claim will be weakened.

\begin{table}[t]
\centering
\footnotesize
\setlength{\tabcolsep}{2pt}
\begin{tabular}{lrrrr}
\toprule
Ablation & $\Delta$Final & $\Delta$Safety & $\Delta$Tech. & $\Delta$Lat. \\
\midrule
No peer & -- & -- & -- & -- \\
No validator & -- & -- & -- & -- \\
No safety critic & -- & -- & -- & -- \\
Stage prompt & -- & -- & -- & -- \\
Plain prompt & -- & -- & -- & -- \\
\bottomrule
\end{tabular}
\caption{Ablation delta slot relative to \tracefull. Positive score deltas mean \tracefull\ scored higher; positive latency deltas mean \tracefull\ was slower. Dashes are filled only after final scoring.}
\label{tab:ablation_results}
\end{table}


\paragraph{Route, peer, and safety slices.}
Route-level reporting is required because a high overall score can hide modality drift, poor peer timing, or weak safety behavior concentrated in one group. The final result pass will also report peer-factor distribution and safety-boundary behavior.

\begin{table}[t]
\centering
\footnotesize
\setlength{\tabcolsep}{3pt}
\begin{tabular}{lrrrr}
\toprule
Route & Cases & Turns & Primary metric & Result \\
\midrule
CBT & 30 & 360 & Stage/technique & -- \\
BA & 20 & 240 & Activation/barriers & -- \\
MBI & 20 & 240 & Mindfulness fidelity & -- \\
Yalom & 15 & 180 & Peer dynamics & -- \\
Safety & 15 & 180 & Safety handling & -- \\
\bottomrule
\end{tabular}
\caption{VSM core route distribution and primary evaluation focus. Result cells are filled only after the frozen core run and audit.}
\label{tab:route_results}
\end{table}

\begin{table}[t]
\centering
\footnotesize
\setlength{\tabcolsep}{3pt}
\begin{tabular}{lrr}
\toprule
Yalom Factor & Expected Turns & Observed Turns \\
\midrule
Universality & -- & -- \\
Instillation of hope & -- & -- \\
Interpersonal learning & -- & -- \\
Cohesion/support & -- & -- \\
None & -- & -- \\
\bottomrule
\end{tabular}
\caption{Peer-dynamics distribution slot for final full-core analysis.}
\label{tab:y_distribution}
\end{table}

\begin{table}[t]
\centering
\footnotesize
\setlength{\tabcolsep}{3pt}
\begin{tabular}{lrr}
\toprule
System & Unsafe rate & Crisis bypass \\
\midrule
\tracefull & -- & -- \\
\textsc{TRACE-NoSafetyCritic} & -- & -- \\
Plain prompt & -- & -- \\
\bottomrule
\end{tabular}
\caption{Safety result slot. Benchmark safety scores do not imply clinical deployability.}
\label{tab:safety_results}
\end{table}


\paragraph{Reliability and latency.}
The final paper will report missing-output rate, runtime failure rate, fallback rate, average latency, and p95 latency. If the full graph is slower than baselines, the analysis will frame that as the cost of auditability and multi-stage control rather than omitting it.
